\documentclass[paper=a4,fontsize=12pt,parskip=false]{scrartcl}
\usepackage[top=14mm,bottom=16mm,left=16mm,right=16mm,headheight=14pt,headsep=4mm,footskip=8mm]{geometry}
\usepackage{fontspec,xcolor,fancyhdr,needspace}
\setmainfont{NotoSerifBengali}[Path=fonts/,Extension=.ttf,UprightFont=*-Regular,BoldFont=*-Bold,Script=Bengali]
\newfontfamily\cmu{cmun}[Path=fonts/,Extension=.ttf,UprightFont=*rm,BoldFont=*bx,ItalicFont=*ti,Scale=1.12]
\newfontfamily\bnsemi{NotoSerifBengali-SemiBold}[Path=fonts/,Extension=.ttf,Script=Bengali]
\newcommand\en[1]{{\cmu #1}}
\definecolor{ink}{HTML}{111111}\definecolor{mute}{HTML}{6E6A62}\definecolor{rule}{HTML}{D5CFC3}
\color{ink}
\linespread{1.32}\selectfont
\setlength\parindent{0pt}\setlength\parskip{0pt}
\raggedright
\pagestyle{fancy}\fancyhf{}
\renewcommand\headrulewidth{0pt}
\fancyhead[L]{\scriptsize\color{mute}\cmu Binomial Theorem \,·\, Narration Script}
\fancyhead[R]{\scriptsize\color{mute}\leftmark}
\fancyfoot[C]{\scriptsize\color{mute}\cmu\thepage}
\newcommand\npart[4]{\vspace{5mm}
  \markboth{\cmu Part #1}{}
  {\large\bfseries #3}\hspace{.8em}{\small\color{mute}\cmu Part #1\normalfont\,·\,#2\,·\,#4টি ক্লিপ}\par
  \vspace{.6mm}{\color{ink}\rule{\linewidth}{1pt}}\par\nobreak\vspace{1mm}\nobreak}
\newcommand\clip[2]{\vspace{1.6mm}
  \noindent{\color{mute}\raisebox{.1ex}{\framebox(7,7){}}}\hspace{.5em}{\bnsemi\small ক্লিপ #1}\hspace{.6em}{\scriptsize\color{mute}\cmu #2}\par\nobreak
  \vspace{.3mm}\nobreak}
\newcommand\sline[1]{\clubpenalty=0\hangindent=1.2em\hspace*{1.2em}\ignorespaces #1\par\vspace{.6mm}}
\newcommand\clipend{\vspace{.8mm}{\color{rule}\rule{\linewidth}{.4pt}}\par}
\begin{document}
\thispagestyle{fancy}
{\small\color{mute}\cmu\addfontfeatures{LetterSpace=10}NARRATION SCRIPT}\par
{\fontsize{20}{26}\selectfont\cmu\bfseries Binomial Theorem}\hspace{.6em}{\large — সংখ্যাগুলো আসলে কোথা থেকে আসে?}\par
{\scriptsize\color{mute} প্রতিটি লাইন একটি বাক্য \,·\, ``—'' চিহ্নে একটু থামুন \,·\, কোড (যেমন \cmu p2\_05\normalfont) = অডিও ফাইলের নাম\par}
\vspace{-2mm}
\npart{1}{\en{Intro}}{ভূমিকা}{৫}
\clip{০১}{p1\_s1a}
\sline{ধরো, আমাদের কাছে দুইটা সংখ্যা আছে — এ আর বি।}
\sline{এখন আমরা যদি এদের যোগ করে \en{power} করি —}
\sline{যেমন ধরো পাওয়ার এক হলে, এ প্লাস বি হোল টু দ্য পাওয়ার ওয়ান, সমান এ প্লাস বি।}
\sline{খুবই সহজ।}
\sline{আবার, পাওয়ার দুই হলে — এ প্লাস বি হোল স্কয়ার, সমান এ স্কয়ার প্লাস টু এ বি প্লাস বি স্কয়ার।}
\sline{এটাও সম্ভব।}
\clipend
\clip{০২}{p1\_s1b}
\sline{আবার পাওয়ার তিন হলে — এ প্লাস বি হোল কিউব, সমান এ কিউব প্লাস থ্রি এ স্কয়ার বি প্লাস থ্রি এ বি স্কয়ার প্লাস বি কিউব।}
\sline{একটু বড়।}
\sline{আর পাওয়ার চার হলে — এ প্লাস বি হোল টু দ্য পাওয়ার ফোর, সমান এ টু দ্য পাওয়ার ফোর প্লাস ফোর এ কিউব বি প্লাস সিক্স এ স্কয়ার বি স্কয়ার প্লাস ফোর এ বি কিউব প্লাস বি টু দ্য পাওয়ার ফোর।}
\sline{একটু কমপ্লিকেটেড, তবে অসম্ভব নয়।}
\clipend
\clip{০৩}{p1\_s2}
\sline{কিন্তু \en{power} যখন আরও বাড়তে থাকবে, তখন কী হবে?}
\sline{যেমন — এ প্লাস বি হোল টু দ্য পাওয়ার ফাইভ, এ প্লাস বি হোল টু দ্য পাওয়ার সিক্স, এ প্লাস বি হোল টু দ্য পাওয়ার সেভেন।}
\sline{তখন \en{expression}-টা ধীরে ধীরে বড় হতে থাকে।}
\sline{আর আমরা যদি প্রতিবার এভাবেই \en{expand} করতে থাকি, তাহলে একটা স্বাভাবিক প্রশ্ন আসে —}
\sline{এর কি কোনো সহজ উপায় নেই?}
\clipend
\clip{০৪}{p1\_s3}
\sline{কারণ একটু ভালো করে লক্ষ্য করলে দেখা যায়, এই \en{expansion}-গুলোর মধ্যে কিছু সংখ্যা বারবার ফিরে আসছে।}
\sline{যেমন, এক।}
\sline{তারপর, এক, এক।}
\sline{তারপর, এক, দুই, এক।}
\sline{তারপর, এক, তিন, তিন, এক।}
\sline{আর এরপর, এক, চার, ছয়, চার, এক।}
\sline{এগুলো কি শুধু \en{coincidence}?}
\sline{নাকি এর পেছনে সত্যিই কোনো \en{pattern} আছে?}
\sline{আর যদি \en{pattern} থেকেই থাকে — তাহলে এই সংখ্যাগুলো আসলে কোথা থেকে আসে?}
\sline{বিশেষ করে এই ছয়-টা।}
\sline{কেন এখানে ছয়?}
\sline{কেন চার নয়?}
\clipend
\clip{০৫}{p1\_s4}
\sline{এই সংখ্যাগুলোকে যদি আমরা শুধু \en{algebra} দিয়ে না দেখে, বরং অন্য একটা দৃষ্টিকোণ থেকে দেখি —}
\sline{তাহলে হয়তো পুরো ব্যাপারটাই অনেক বেশি সহজ হয়ে যাবে।}
\sline{আর সেই অন্য দৃষ্টিকোণটা হলো —}
\sline{\en{combination}।}
\clipend
\npart{2}{অধ্যায় ১}{\en{Combination}}{১৭}
\clip{০১}{p2\_01}
\sline{\en{Combination}-এর পুরো বিষয়টা বুঝতে গেলে, প্রথমেই আমাদের একটা গুরুত্বপূর্ণ জিনিস বুঝতে হবে — \en{factorial}।}
\sline{সাধারণভাবে আমরা জানি, এন ফ্যাক্টোরিয়াল সমান ওয়ান ইনটু টু ইনটু থ্রি ইনটু ফোর ইনটু, এভাবে, ইনটু এন।}
\sline{কিন্তু একটা প্রশ্ন হলো — এই এন ফ্যাক্টোরিয়াল আসলে আমাদের কী বোঝায়?}
\sline{শুধু এক থেকে এন পর্যন্ত সংখ্যাগুলো গুণ করলেই তো হলো না।}
\sline{এর পেছনে আসলে একটা খুব সুন্দর \en{idea} আছে।}
\sline{\en{factorial} মূলত বলে — এন টি আলাদা বস্তুকে কতভাবে সাজানো যায়, যেখানে তাদের \en{position} বা \en{order} গুরুত্বপূর্ণ।}
\clipend
\clip{০২}{p2\_02}
\sline{এটা বুঝতে আমরা একটা খুব ছোট উদাহরণ দেখি।}
\sline{ধরো, আমাদের কাছে তিনজন মানুষ আছে — এ, বি, সি।}
\sline{এখন এই তিনজনকে একটা লাইনে সাজাতে হবে।}
\sline{খেয়াল করো, এখানে মানুষ তিনজনই একই থাকছে।}
\sline{কেউ দল থেকে বেরিয়ে যাচ্ছে না, নতুন কেউও আসছে না।}
\sline{আমরা শুধু দেখছি — এই তিনজনকে বিভিন্ন \en{position}-এ কতভাবে বসানো যায়।}
\clipend
\clip{০৩}{p2\_03}
\sline{অনেকটা একটা ফুটবল দলের মতো।}
\sline{দলের খেলোয়াড় একই থাকতে পারে, কিন্তু তাদের \en{position} পরিবর্তন করলে \en{formation} পুরো বদলে যায়।}
\sline{যেমন — যে গোলকিপার, সে যদি স্ট্রাইকারে খেলে, তাহলে দলের মান কিন্তু আগের মতো থাকবে না।}
\sline{আবার যে স্ট্রাইকার, সে যদি ডিফেন্সে খেলে, তাহলে দলটাকে পুরো ভিন্ন মনে হবে।}
\sline{ঠিক একইভাবে এখানে — এ, বি, সি — এই তিনজনই থাকবে।}
\sline{কিন্তু তাদের \en{position} পরিবর্তন করে আমরা কতগুলো আলাদা \en{arrangement} তৈরি করতে পারি, সেটাই আমাদের জানতে হবে।}
\clipend
\clip{০৪}{p2\_04}
\sline{এখন চলো দেখি, এটা কীভাবে গোনা যায়।}
\sline{আমাদের মোট তিনটা \en{position} আছে।}
\sline{প্রথম \en{position}-এ কে বসবে?}
\sline{এখানে এ, বি, অথবা সি — যেকোনো একজন বসতে পারে।}
\sline{অর্থাৎ প্রথম \en{position}-এর জন্য আমাদের কাছে আছে — তিনটি \en{choice}।}
\sline{ধরো, প্রথম \en{position}-এ আমরা এ-কে বসালাম।}
\sline{এখন দ্বিতীয় \en{position}-এ এ আর বসতে পারবে না, কারণ এ ইতিমধ্যেই বসে গেছে।}
\sline{তাই দ্বিতীয় \en{position}-এ বাকি আছে — বি, সি।}
\sline{অর্থাৎ এখানে আছে — দুইটি \en{choice}।}
\clipend
\clip{০৫}{p2\_05}
\sline{এখন ধরো, দ্বিতীয় \en{position}-এ বি বসলো।}
\sline{তাহলে তৃতীয় \en{position}-এর জন্য আর বাকি থাকল শুধু সি।}
\sline{অর্থাৎ শেষ \en{position}-এর জন্য আছে — একটি \en{choice}।}
\sline{তাহলে মোট \en{arrangement}-এর সংখ্যা হবে — থ্রি ইনটু টু ইনটু ওয়ান।}
\sline{অর্থাৎ, ছয়।}
\clipend
\clip{০৬}{p2\_06}
\sline{কিন্তু এখানে একটা গুরুত্বপূর্ণ ব্যাপার আছে।}
\sline{আমরা শুধু একটা \en{arrangement} পাইনি।}
\sline{প্রথম \en{position}-এ এ বসালে, দ্বিতীয় \en{position}-এ বি অথবা সি — দুইভাবেই যেতে পারি।}
\sline{আবার প্রথম \en{position}-এ বি বসালেও, পরের \en{position}-এর জন্য দুইটা \en{choice} থাকে।}
\sline{একইভাবে প্রথম \en{position}-এ সি বসালেও দুইটা \en{choice} থাকে।}
\sline{অর্থাৎ প্রথমে আমাদের ছিল তিনটা পথ, প্রতিটি পথের আবার ছিল দুইটা করে পথ, আর প্রতিটি সেই পথের শেষে ছিল একটা করে পথ।}
\sline{তাই — থ্রি ইনটু টু ইনটু ওয়ান সমান সিক্স।}
\sline{এবং ঠিক এ কারণেই আমরা ছয়টি \en{arrangement} পাই: এ বি সি, এ সি বি, বি এ সি, বি সি এ, সি এ বি, সি বি এ।}
\clipend
\clip{০৭}{p2\_07}
\sline{এখন এই একই চিন্তাটাকে যদি তিনজনের বদলে এন জন মানুষের ক্ষেত্রে ব্যবহার করি —}
\sline{যেমন, প্রথম \en{position}-এর জন্য থাকবে: এন টি \en{choice}।}
\sline{তাহলে, একজন বসে যাওয়ার পর দ্বিতীয় \en{position}-এর জন্য থাকবে: এন মাইনাস ওয়ান টি \en{choice}।}
\sline{তারপর: এন মাইনাস টু।}
\sline{তারপর: এন মাইনাস থ্রি।}
\sline{এভাবে চলতে চলতে শেষ পর্যন্ত একে গিয়ে থামবে।}
\sline{তাই মোট \en{arrangement} হবে — এন ইনটু এন মাইনাস ওয়ান ইনটু এন মাইনাস টু ইনটু এন মাইনাস থ্রি, এভাবে, ইনটু টু ইনটু ওয়ান।}
\sline{এই পুরো \en{multiplication}-টাকে সংক্ষেপে আমরা লিখি — এন ফ্যাক্টোরিয়াল।}
\sline{এটা আমাদের বলে — এন টি আলাদা বস্তুকে \en{order} গুরুত্বপূর্ণ রেখে কতভাবে সাজানো যায়।}
\clipend
\clip{০৮}{p2\_08}
\sline{আবার, এখানেই আমাদের পরের প্রশ্নটা শুরু হয় —}
\sline{যদি \en{order} আর গুরুত্বপূর্ণ না হয়, তখন কী হবে?}
\sline{ধরো, আমাদের কাছে পাঁচজন মানুষ আছে — এ, বি, সি, ডি, ই।}
\sline{এখন আমরা যদি এদের মধ্যে থেকে দুইজন করে দল তৈরি করতে চাই, তাহলে মোট কতগুলো আলাদা দল তৈরি করা সম্ভব?}
\sline{তো এইটা সমাধান করার জন্য আমরা এভাবে চিন্তা করতে পারি যে, প্রথম \en{position}-এর জন্য পাঁচজনের যেকোনো একজনকে নিতে পারি।}
\sline{তাই আছে — পাঁচটি \en{choice}।}
\sline{একজনকে নেওয়ার পর বাকি থাকে চারজন।}
\sline{তাই দ্বিতীয় \en{position}-এর জন্য আছে — চারটি \en{choice}।}
\sline{সুতরাং মোট \en{ordered selection} হবে — ফাইভ ইনটু ফোর সমান টুয়েন্টি।}
\clipend
\clip{০৯}{p2\_09}
\sline{এবার যদি আমরা তিনজন করে দল তৈরি করতে চাই?}
\sline{প্রথম \en{position}-এর জন্য থাকবে পাঁচটি \en{choice}।}
\sline{একজন নেওয়ার পর দ্বিতীয় \en{position}-এর জন্য থাকবে চারটি।}
\sline{আর দুজন নেওয়ার পর তৃতীয় \en{position}-এর জন্য থাকবে তিনটি।}
\sline{তাই — ফাইভ ইনটু ফোর ইনটু থ্রি সমান সিক্সটি টি \en{ordered selection}।}
\sline{খেয়াল করো, এখানে একটা \en{pattern} তৈরি হচ্ছে।}
\sline{দুইজন নিলে: ফাইভ ইনটু ফোর।}
\sline{তিনজন নিলে: ফাইভ ইনটু ফোর ইনটু থ্রি।}
\sline{চারজন নিলে: ফাইভ ইনটু ফোর ইনটু থ্রি ইনটু টু।}
\sline{অর্থাৎ প্রতিবার আমরা যখন একজন করে মানুষ নিই, পরের \en{position}-এর জন্য \en{available} মানুষের সংখ্যা এক করে কমে যায়।}
\clipend
\clip{১০}{p2\_10}
\sline{তাহলে যদি আমাদের কাছে এন টি বস্তু থাকে এবং আমরা এম টি বস্তু নিতে চাই, তাহলে প্রথম \en{position}-এর জন্য থাকবে — এন টি \en{choice}।}
\sline{দ্বিতীয় \en{position}-এর জন্য — এন মাইনাস ওয়ান টি।}
\sline{তৃতীয় \en{position}-এর জন্য — এন মাইনাস টু টি।}
\sline{এভাবে চলতে চলতে এম-তম \en{position}-এ পৌঁছালে, তার আগে আমরা এম মাইনাস ওয়ান টি বস্তু নিয়ে ফেলেছি।}
\sline{তাই তখন বাকি থাকবে — এন মাইনাস, ব্র্যাকেটে এম মাইনাস ওয়ান।}
\sline{অর্থাৎ, এন মাইনাস এম প্লাস ওয়ান টি বস্তু।}
\sline{তাই \en{order} গুরুত্বপূর্ণ রেখে এম টি বস্তু নেওয়ার মোট উপায় হবে — এন ইনটু এন মাইনাস ওয়ান ইনটু এন মাইনাস টু, এভাবে, ইনটু এন মাইনাস এম প্লাস ওয়ান।}
\clipend
\clip{১১}{p2\_11}
\sline{কিন্তু এখানে একটা সমস্যা আছে।}
\sline{দেখো, আমরা একই দলকে দুইবার করে গুনে ফেলছি — শুধুমাত্র তাদের \en{position} পরিবর্তন হওয়ার কারণে।}
\sline{যেমন, এ এবং বি-কে নিলে আমরা একবার এ বি পাচ্ছি, আবার বি এ পাচ্ছি।}
\sline{কিন্তু এ বি আর বি এ-তে মানুষের সংখ্যা বা দলের সদস্য কেউই বদলায়নি।}
\sline{শুধু তাদের \en{position} পরিবর্তন হয়েছে।}
\sline{অথচ আমাদের উদ্দেশ্য হলো কতগুলো আলাদা দল তৈরি করা যায়, সেটা বের করা — কতভাবে একই দলকে সাজানো যায়, সেটা নয়।}
\sline{তাই এখন আমাদের বের করতে হবে — একই দলকে আমরা মোট কতবার করে গুনে ফেলেছি।}
\clipend
\clip{১২}{p2\_12}
\sline{এখন এই সমস্যাটা একটু ছোট করে দেখি।}
\sline{ধরো, আমাদের কাছে শুধু দুইজন মানুষ আছে — এ, বি।}
\sline{এদের দুজনকে নিয়ে দল বানাতে চাই।}
\sline{\en{Order} গুরুত্বপূর্ণ ধরে আমরা লিখতে পারি: এ বি।}
\sline{আবার লিখতে পারি: বি এ।}
\sline{কিন্তু এ বি আর বি এ কি সত্যিই দুইটা আলাদা দল?}
\sline{না।}
\sline{দুটোতেই তো একই দুইজন মানুষ আছে — এ, বি।}
\sline{শুধু তাদের \en{position} পরিবর্তন হয়েছে।}
\sline{অর্থাৎ আমরা একই দলকে দুইবার গুনে ফেলেছি।}
\sline{আর দুইজন বস্তুকে \en{order} গুরুত্বপূর্ণ রেখে সাজানো যায় — টু ফ্যাক্টোরিয়াল সমান টু ভাবে।}
\sline{তাই একই দল আমাদের গণনার মধ্যে দুইবার এসেছে।}
\clipend
\clip{১৩}{p2\_13}
\sline{এবার তিনজনের ক্ষেত্রে দেখি।}
\sline{ধরো, এ, বি, সি — এই তিনজনকে নিয়ে দল তৈরি করছি।}
\sline{\en{Order} গুরুত্বপূর্ণ ধরে তাদের সাজালে পাই — এ বি সি, এ সি বি, বি এ সি, বি সি এ, সি এ বি, সি বি এ।}
\sline{অর্থাৎ মোট — থ্রি ইনটু টু ইনটু ওয়ান সমান সিক্স টি \en{arrangement}।}
\sline{কিন্তু এখানে খেয়াল করো — এই ছয়টি \en{arrangement}-এর মধ্যে দল একবারও পরিবর্তন হয়নি।}
\sline{প্রতিবারই দলটি একই: এ, বি, সি।}
\sline{শুধু \en{position} পরিবর্তন হয়েছে।}
\sline{তাই একই দলকে আমরা মোট ছয় বার গুনেছি।}
\sline{আর থ্রি ফ্যাক্টোরিয়াল সমান সিক্স।}
\sline{তাই মোট \en{arrangement}-কে থ্রি ফ্যাক্টোরিয়াল দিয়ে ভাগ করলে আসল দলের সংখ্যা পাওয়া যাবে: সিক্স বাই থ্রি ফ্যাক্টোরিয়াল, সমান সিক্স বাই সিক্স, সমান ওয়ান।}
\sline{অর্থাৎ তিনজন মানুষ থেকে তিনজনের দল আসলে — একটি।}
\clipend
\clip{১৪}{p2\_14}
\sline{এখন দুইজনের ক্ষেত্রে আমরা দেখলাম একই দল এসেছে টু ফ্যাক্টোরিয়াল বার।}
\sline{তিনজনের ক্ষেত্রে এসেছে থ্রি ফ্যাক্টোরিয়াল বার।}
\sline{তাহলে সাধারণভাবে যদি একটি দলে এম টি বস্তু থাকে, সেই এম টি বস্তুকে \en{order} গুরুত্বপূর্ণ রেখে সাজানো যাবে — এম ফ্যাক্টোরিয়াল ভাবে।}
\sline{অর্থাৎ আমাদের আগের গণনায় প্রতিটি আসল দল এম ফ্যাক্টোরিয়াল বার করে গোনা হয়েছে।}
\sline{আর আমরা যেহেতু \en{order} চাই না, তাই এই অতিরিক্ত গণনাগুলো সরিয়ে দিতে এম ফ্যাক্টোরিয়াল দিয়ে ভাগ করতে হবে।}
\clipend
\clip{১৫}{p2\_15}
\sline{তাহলে আমাদের পুরো \en{equation} দাঁড়ায় — এন ইনটু এন মাইনাস ওয়ান ইনটু এন মাইনাস টু, এভাবে, ইনটু এন মাইনাস এম প্লাস ওয়ান, বাই এম ফ্যাক্টোরিয়াল।}
\sline{আবার আরো সাধারণ রূপে লিখলে —}
\sline{সমান, এন ইনটু এন মাইনাস ওয়ান, এভাবে, ইনটু এন মাইনাস এম প্লাস ওয়ান, ইনটু এন মাইনাস এম ফ্যাক্টোরিয়াল, বাই এম ফ্যাক্টোরিয়াল ইনটু এন মাইনাস এম ফ্যাক্টোরিয়াল।}
\sline{সমান, এন ফ্যাক্টোরিয়াল বাই এম ফ্যাক্টোরিয়াল ইনটু এন মাইনাস এম ফ্যাক্টোরিয়াল।}
\sline{এটাই এন টি বস্তুর মধ্যে থেকে এম টি বস্তু বেছে নিয়ে কতগুলো আলাদা দল তৈরি করা যায়, তার সংখ্যা।}
\sline{আর এই সংখ্যাকেই আমরা লিখি — সি এন এম।}
\sline{আর এটাই হলো \en{Combination}-এর অত্যন্ত গুরুত্বপূর্ণ একটি ধারণা।}
\clipend
\clip{১৬}{p2\_16}
\sline{এখন লক্ষ্য করো, এ প্লাস বি হোল টু দ্য পাওয়ার ওয়ান, সমান এ প্লাস বি।}
\sline{আবার, এ প্লাস বি হোল স্কয়ার, সমান এ স্কয়ার প্লাস টু এ বি প্লাস বি স্কয়ার।}
\sline{এবং, এ প্লাস বি হোল কিউব, সমান এ কিউব প্লাস থ্রি এ স্কয়ার বি প্লাস থ্রি এ বি স্কয়ার প্লাস বি কিউব।}
\sline{এভাবেই \en{power} বাড়তে থাকলে আমরা আরও বড় বড় \en{expansion} পেতে থাকি।}
\sline{কিন্তু এখানে একটা বিষয় লক্ষ্য করো।}
\sline{প্রতিটি \en{expansion}-এর \en{term}-এর সামনে কিছু সংখ্যা আছে।}
\sline{যেমন, এক।}
\sline{আবার, এক, এক।}
\sline{আবার, এক, দুই, এক।}
\sline{আবার, এক, তিন, তিন, এক।}
\sline{এই সংখ্যাগুলোই হলো আমাদের \en{coefficient}, বা সহগ।}
\clipend
\clip{১৭}{p2\_17}
\sline{এই \en{coefficient}-গুলো আমরা কীভাবে বের করব?}
\sline{কারণ প্রতিটি \en{term}-এর সাথে প্রতিটি \en{term}-এর গুণফল করে শেষ পর্যায়ে এই সহগগুলো বের করার প্রচেষ্টা আসলে খুবই দীর্ঘ।}
\sline{তাই প্রশ্ন আসে — যদি একটা \en{pattern} খুঁজে বের করা যায় এবং সেটা অনুযায়ী এই \en{coefficient}-গুলো সরাসরি বের করা যায়, তাহলে?}
\sline{আর এখান থেকেই সামনে আসে নতুন একটা \en{concept} —}
\sline{\en{Binomial Theorem}।}
\clipend
\npart{3}{অধ্যায় ২}{\en{Binomial Theorem}}{১৭}
\clip{০১}{p3\_01}
\sline{এই ব্যাপারটা বোঝার জন্য আমরা আগে খুব ছোট একটা উদাহরণ দেখি।}
\sline{ধরো, এ প্লাস বি, ইনটু, এ প্লাস বি।}
\sline{অর্থাৎ, এ প্লাস বি হোল স্কয়ার।}
\sline{এখন সরাসরি এ স্কয়ার প্লাস টু এ বি প্লাস বি স্কয়ার না লিখে, আমরা এটাকে একটু অন্যভাবে দেখি।}
\sline{বোঝার সুবিধার জন্য প্রথম এ প্লাস বি-কে আমরা লিখি — এ ওয়ান প্লাস বি ওয়ান।}
\sline{এবং দ্বিতীয় এ প্লাস বি-কে লিখি — এ টু প্লাস বি টু।}
\sline{তাহলে পুরো \en{expression} হবে — এ ওয়ান প্লাস বি ওয়ান, ইনটু, এ টু প্লাস বি টু।}
\clipend
\clip{০২}{p3\_02}
\sline{এখানে এ ওয়ান আর এ টু একই এ-এর দুইটি আলাদা \en{label}।}
\sline{অর্থাৎ এরা একই জিনিসকে বোঝাচ্ছে, কিন্তু আমরা \en{label} আলাদা করেছি শুধু এটা বোঝার জন্য যে প্রথম \en{factor} থেকে কোন এ এসেছে আর দ্বিতীয় \en{factor} থেকে কোন এ এসেছে।}
\sline{একইভাবে, বি ওয়ান, বি টু হলো বি-এর দুইটি আলাদা \en{label}।}
\clipend
\clip{০৩}{p3\_03}
\sline{এখন আমরা \en{expression}-টা \en{expand} করি।}
\sline{এ ওয়ান প্লাস বি ওয়ান, ইনটু, এ টু প্লাস বি টু হলে পাই — এ ওয়ান এ টু, প্লাস এ ওয়ান বি টু, প্লাস বি ওয়ান এ টু, প্লাস বি ওয়ান বি টু।}
\sline{এখন একটু থেমে প্রতিটি \en{term} দেখি।}
\sline{প্রথমে, এ ওয়ান এ টু।}
\sline{এখানে বি-এর কোনো \en{choice} নেই।}
\sline{অর্থাৎ দুইটি \en{factor}-এর দুটিতেই আমরা এ নিয়েছি।}
\sline{তারপর, বি ওয়ান বি টু।}
\sline{এখানে ঠিক উল্টো — দুইটি \en{factor} থেকেই বি নেওয়া হয়েছে।}
\clipend
\clip{০৪}{p3\_04}
\sline{কিন্তু মাঝখানে দুটি \en{term} আছে: এ ওয়ান বি টু, এবং বি ওয়ান এ টু।}
\sline{এখানে কী ঘটছে?}
\sline{দুইটি \en{factor}-এর মধ্যে একটিতে বি এবং অন্যটিতে এ এসেছে।}
\sline{শুধু বি কোন \en{factor} থেকে এসেছে, সেটা বদলেছে।}
\sline{প্রথমটিতে বি এসেছে দ্বিতীয় \en{factor} থেকে: এ ওয়ান বি টু।}
\sline{আর দ্বিতীয়টিতে বি এসেছে প্রথম \en{factor} থেকে: বি ওয়ান এ টু।}
\sline{অর্থাৎ বি নেওয়ার জন্য আমাদের কাছে দুইটি \en{possibility} আছে।}
\clipend
\clip{০৫}{p3\_05}
\sline{তাই এই দুইটি \en{term} পরে একই \en{algebraic term} হয়ে যাবে: এ ওয়ান বি টু প্লাস বি ওয়ান এ টু, সমান এ বি প্লাস এ বি, সমান টু এ বি।}
\sline{এখন আমরা পুরো \en{expression}-টাকে আবার লিখতে পারি — এ ওয়ান এ টু, প্লাস এ ওয়ান বি টু, প্লাস বি ওয়ান এ টু, প্লাস বি ওয়ান বি টু।}
\sline{অর্থাৎ, এ স্কয়ার প্লাস টু এ বি প্লাস বি স্কয়ার।}
\sline{এখানে যে টু এসেছে কারণ দুইটি \en{factor}-এর মধ্যে যেকোনো একটি \en{factor} থেকে বি নেওয়া যায়।}
\clipend
\clip{০৬}{p3\_06}
\sline{এখন আরেকটা উদাহরণ দেওয়া যাক।}
\sline{ধরো, এ প্লাস বি হোল কিউব।}
\sline{এটাকে লিখলাম — এ ওয়ান প্লাস বি ওয়ান, এ টু প্লাস বি টু, এ থ্রি প্লাস বি থ্রি।}
\sline{এখন তিনটি \en{factor} থেকে একটি করে \en{choice} নিয়ে \en{expansion} করলে আমরা পাই —}
\sline{এ ওয়ান এ টু এ থ্রি, প্লাস এ ওয়ান এ টু বি থ্রি, প্লাস এ ওয়ান বি টু এ থ্রি, প্লাস বি ওয়ান এ টু এ থ্রি,}
\sline{প্লাস এ ওয়ান বি টু বি থ্রি, প্লাস বি ওয়ান এ টু বি থ্রি, প্লাস বি ওয়ান বি টু এ থ্রি,}
\sline{প্লাস বি ওয়ান বি টু বি থ্রি।}
\sline{এখন এটাকে আমরা একটু অন্যভাবে সাজাই।}
\clipend
\clip{০৭}{p3\_07}
\sline{প্রথম \en{term}: এ ওয়ান এ টু এ থ্রি।}
\sline{এখানে বি-এর কোনো \en{choice} নেই।}
\sline{অর্থাৎ বি নেওয়া হয়েছে শূন্য বার।}
\sline{তারপর তিনটি \en{term}: এ ওয়ান এ টু বি থ্রি, এ ওয়ান বি টু এ থ্রি, বি ওয়ান এ টু এ থ্রি।}
\sline{এখানে প্রতিটি \en{term}-এ বি এসেছে ঠিক এক বার।}
\sline{তারপর আবার তিনটি \en{term}: এ ওয়ান বি টু বি থ্রি, বি ওয়ান এ টু বি থ্রি, বি ওয়ান বি টু এ থ্রি।}
\sline{এখানে বি এসেছে ঠিক দুই বার।}
\sline{আর একদম শেষে: বি ওয়ান বি টু বি থ্রি।}
\sline{এখানে বি এসেছে তিন বার।}
\clipend
\clip{০৮}{p3\_08}
\sline{এখন একটা সুন্দর \en{pattern} দেখা যাচ্ছে।}
\sline{জিরো বি, ওয়ান বি, টু বি, থ্রি বি।}
\sline{অর্থাৎ আমরা \en{expansion}-টাকে বি-এর সংখ্যা অনুযায়ী ভাগ করতে পারি।}
\sline{এখন প্রথমটার দিকে তাকাই।}
\sline{এ ওয়ান এ টু এ থ্রি।}
\sline{এখানে তিনটি \en{factor}-এর কোনোটিতেই বি আসেনি।}
\sline{অর্থাৎ বি-এর তিনটি \en{possible position}-এর মধ্যে থেকে শূন্যটি \en{position} বেছে নেওয়া হয়েছে।}
\sline{এটা \en{Combination}-এর ভাষায়: সি থ্রি জিরো সমান ওয়ান।}
\sline{তাই এর \en{coefficient} হবে ওয়ান।}
\clipend
\clip{০৯}{p3\_09}
\sline{এখন পরের তিনটি \en{term} দেখি।}
\sline{এ ওয়ান এ টু বি থ্রি, এ ওয়ান বি টু এ থ্রি, বি ওয়ান এ টু এ থ্রি।}
\sline{এখানে বি এসেছে একবার।}
\sline{তাহলে প্রশ্নটা দাঁড়াল — তিনটি \en{factor}-এর মধ্যে কোন একটি \en{factor} থেকে বি নেব?}
\sline{প্রথম \en{factor} থেকে নিতে পারি।}
\sline{দ্বিতীয় \en{factor} থেকে নিতে পারি।}
\sline{তৃতীয় \en{factor} থেকেও নিতে পারি।}
\sline{অর্থাৎ মোট তিনটি উপায়।}
\sline{\en{Combination} দিয়ে, সি থ্রি ওয়ান সমান থ্রি।}
\sline{তাই এই তিনটি \en{term} একসাথে হয়ে যাবে — থ্রি এ স্কয়ার বি।}
\clipend
\clip{১০}{p3\_10}
\sline{এবার পরের তিনটি \en{term} দেখি।}
\sline{এ ওয়ান বি টু বি থ্রি, বি ওয়ান এ টু বি থ্রি, বি ওয়ান বি টু এ থ্রি।}
\sline{এখানে বি এসেছে দুইবার।}
\sline{তাহলে এবার প্রশ্ন — তিনটি \en{factor}-এর মধ্যে থেকে কোন দুইটি \en{factor}-এ বি থাকবে?}
\sline{প্রথম এবং দ্বিতীয়, প্রথম এবং তৃতীয়, অথবা দ্বিতীয় এবং তৃতীয়।}
\sline{অর্থাৎ মোট তিনটি উপায়।}
\sline{তাই, সি থ্রি টু সমান থ্রি।}
\sline{এবং এই তিনটি \en{term} মিলে হবে — থ্রি এ বি স্কয়ার।}
\clipend
\clip{১১}{p3\_11}
\sline{আর একদম শেষে: বি ওয়ান বি টু বি থ্রি।}
\sline{এখানে তিনটি \en{factor}-এর তিনটিতেই বি নেওয়া হয়েছে।}
\sline{অর্থাৎ তিনটি থেকে তিনটিই বেছে নেওয়া হয়েছে।}
\sline{তাই, সি থ্রি থ্রি সমান ওয়ান।}
\sline{এবং \en{term} হলো — বি কিউব।}
\sline{তাহলে পুরো \en{expansion} হয়ে গেল — এ কিউব প্লাস থ্রি এ স্কয়ার বি প্লাস থ্রি এ বি স্কয়ার প্লাস বি কিউব।}
\sline{অথবা, সি থ্রি জিরো এ কিউব, প্লাস সি থ্রি ওয়ান এ স্কয়ার বি, প্লাস সি থ্রি টু এ বি স্কয়ার, প্লাস সি থ্রি থ্রি বি কিউব।}
\clipend
\clip{১২}{p3\_12}
\sline{এখন একটু খেয়াল করো, এখানে এ আর বি-এর \en{power}-এর মধ্যেও একটা সুন্দর নিয়ম আছে।}
\sline{যখন বি-এর \en{power} জিরো, তখন এ-এর \en{power} থ্রি।}
\sline{যখন বি-এর \en{power} ওয়ান, তখন এ-এর \en{power} টু।}
\sline{যখন বি-এর \en{power} টু, তখন এ-এর \en{power} ওয়ান।}
\sline{আর যখন বি-এর \en{power} থ্রি, তখন এ-এর \en{power} জিরো।}
\sline{অর্থাৎ প্রতিবার বি-এর \en{power} এক করে বাড়ছে, আর এ-এর \en{power} এক করে কমছে।}
\sline{কারণ মোট \en{factor} সবসময়ই তিনটি।}
\sline{তাই, এ কিউব, এ স্কয়ার বি, এ বি স্কয়ার, বি কিউব — প্রতিটি \en{term}-এ এ আর বি-এর \en{power} যোগ করলে সবসময় থ্রি পাওয়া যায়।}
\clipend
\clip{১৩}{p3\_13}
\sline{এখন এই একই নিয়ম যদি থ্রি-এর বদলে যেকোনো এন-এর জন্য ব্যবহার করি, তাহলে এন টি \en{factor}-এর মধ্যে থেকে আর টি \en{position} বেছে নিয়ে বি নিতে পারব — সি এন আর ভাবে।}
\sline{আর যেহেতু বাকি এন মাইনাস আর টি \en{position} থেকে এ আসবে, তাই সেই \en{term} হবে — এ টু দ্য পাওয়ার এন মাইনাস আর, বি টু দ্য পাওয়ার আর।}
\sline{অর্থাৎ প্রতিটি \en{term}-এর সাধারণ রূপ: সি এন আর, এ টু দ্য পাওয়ার এন মাইনাস আর, বি টু দ্য পাওয়ার আর।}
\clipend
\clip{১৪}{p3\_14}
\sline{আর সবগুলো \en{term} একসাথে লিখলে আমরা পাই —}
\sline{এ প্লাস বি হোল টু দ্য পাওয়ার এন, সমান সিগমা আর সমান জিরো থেকে এন, সি এন আর, এ টু দ্য পাওয়ার এন মাইনাস আর, বি টু দ্য পাওয়ার আর।}
\sline{এটাই হলো \en{Binomial Theorem}-এর মূল ধারণা।}
\sline{আর এই সংখ্যাটাই হয়ে যাচ্ছে আমাদের \en{coefficient}।}
\clipend
\clip{১৫}{p3\_15}
\sline{এখন চলো, এন সমান টু-এর ক্ষেত্রে একটা উদাহরণ দেখি।}
\sline{আমাদের \en{equation} হবে — এ প্লাস বি হোল স্কয়ার, সমান সিগমা আর সমান জিরো থেকে টু, সি টু আর, এ টু দ্য পাওয়ার টু মাইনাস আর, বি টু দ্য পাওয়ার আর।}
\sline{অর্থাৎ তিনটি \en{term} থাকবে: সি টু জিরো এ স্কয়ার, প্লাস সি টু ওয়ান এ বি, প্লাস সি টু টু বি স্কয়ার।}
\sline{এখন, সি টু জিরো সমান ওয়ান, সি টু ওয়ান সমান টু, সি টু টু সমান ওয়ান।}
\sline{তাই, এ প্লাস বি হোল স্কয়ার সমান এ স্কয়ার প্লাস টু এ বি প্লাস বি স্কয়ার।}
\sline{অর্থাৎ আমরা শুরুতে যে \en{expansion} দেখেছিলাম, আমাদের \en{general rule} ব্যবহার করেও ঠিক একই \en{result} ফিরে পেলাম।}
\clipend
\clip{১৬}{p3\_16}
\sline{এখন \en{general equation}-এর দিকে লক্ষ্য করো।}
\sline{এ প্লাস বি হোল টু দ্য পাওয়ার এন, সমান সিগমা আর সমান জিরো থেকে এন, সি এন আর, এ টু দ্য পাওয়ার এন মাইনাস আর, বি টু দ্য পাওয়ার আর।}
\sline{এখানে আমরা দেখলাম, \en{coefficient}-এর জায়গায় যে সংখ্যাটা আসছে, সেটা আসলে \en{Combination} থেকেই এসেছে।}
\sline{আর \en{Combination}-এর সূত্র যদি বসিয়ে দিই, তাহলে পাই —}
\sline{এ প্লাস বি হোল টু দ্য পাওয়ার এন, সমান সিগমা আর সমান জিরো থেকে এন, এন ফ্যাক্টোরিয়াল বাই আর ফ্যাক্টোরিয়াল ইনটু এন মাইনাস আর ফ্যাক্টোরিয়াল, এ টু দ্য পাওয়ার এন মাইনাস আর, বি টু দ্য পাওয়ার আর।}
\clipend
\clip{১৭}{p3\_17}
\sline{এখন এ ওয়ান, এ টু, এ থ্রি এবং বি ওয়ান, বি টু, বি থ্রি-এর মতো \en{label}-গুলোর আর প্রয়োজন নেই।}
\sline{ওগুলো আমরা শুধু ধারণাটা পরিষ্কার করার জন্য ব্যবহার করেছিলাম।}
\sline{কারণ বাস্তবে একই এ-এর আলাদা আলাদা \en{label} না লিখে আমরা সরাসরি এ টু দ্য পাওয়ার কে লিখতে পারি।}
\sline{আর তখন আমাদের \en{Binomial Theorem}-এর পরিচিত রূপটি সামনে আসে —}
\sline{এ প্লাস বি হোল টু দ্য পাওয়ার এন, সমান সি এন জিরো এ টু দ্য পাওয়ার এন, প্লাস সি এন ওয়ান এ টু দ্য পাওয়ার এন মাইনাস ওয়ান বি, প্লাস সি এন টু এ টু দ্য পাওয়ার এন মাইনাস টু বি স্কয়ার, প্লাস, এভাবে, প্লাস সি এন এন বি টু দ্য পাওয়ার এন।}
\clipend
\npart{4}{অধ্যায় ৩ · সারাংশ}{\en{Pascal's Triangle} ও \en{Summary}}{১৯}
\clip{০১}{p4\_01}
\sline{আমরা এতক্ষণ \en{Binomial Theorem}-এর মাধ্যমে দেখলাম — এ প্লাস বি হোল টু দ্য পাওয়ার এন-কে কীভাবে \en{expand} করা যায়।}
\sline{এবং আমরা এটাও দেখলাম, প্রতিটি \en{term}-এর সামনে যে সংখ্যাগুলো থাকে, সেগুলো আসলে \en{Combination} থেকে আসে।}
\sline{এখন একটু অন্যভাবে বিষয়টাকে দেখি।}
\sline{ধরো আমরা একে একে লিখলাম —}
\sline{এ প্লাস বি হোল টু দ্য পাওয়ার জিরো, সমান ওয়ান।}
\sline{এ প্লাস বি হোল টু দ্য পাওয়ার ওয়ান, সমান এ প্লাস বি।}
\sline{এ প্লাস বি হোল স্কয়ার, সমান এ স্কয়ার প্লাস টু এ বি প্লাস বি স্কয়ার।}
\sline{এ প্লাস বি হোল কিউব, সমান এ কিউব প্লাস থ্রি এ স্কয়ার বি প্লাস থ্রি এ বি স্কয়ার প্লাস বি কিউব।}
\sline{এ প্লাস বি হোল টু দ্য পাওয়ার ফোর, সমান এ টু দ্য পাওয়ার ফোর প্লাস ফোর এ কিউব বি প্লাস সিক্স এ স্কয়ার বি স্কয়ার প্লাস ফোর এ বি কিউব প্লাস বি টু দ্য পাওয়ার ফোর।}
\clipend
\clip{০২}{p4\_02}
\sline{এভাবে আরও কয়েকটি \en{power} লিখলে আমরা পাই — এ প্লাস বি হোল টু দ্য পাওয়ার ফাইভ, এ প্লাস বি হোল টু দ্য পাওয়ার সিক্স, ইত্যাদি।}
\sline{এখন এখানে এ আর বি-এর দিকে না তাকিয়ে, শুধু প্রতিটি \en{term}-এর \en{coefficient}-গুলো আলাদা করে দেখি।}
\sline{প্রথমে পাই — ওয়ান।}
\sline{তারপর — ওয়ান, ওয়ান।}
\sline{তারপর — ওয়ান, টু, ওয়ান।}
\sline{তারপর — ওয়ান, থ্রি, থ্রি, ওয়ান।}
\sline{তারপর — ওয়ান, ফোর, সিক্স, ফোর, ওয়ান।}
\clipend
\clip{০৩}{p4\_03}
\sline{এভাবে \en{coefficient}-গুলোকে যদি আমরা একটার নিচে আরেকটা সাজাই, তাহলে একটা অদ্ভুত সুন্দর \en{pattern} তৈরি হয়।}
\sline{দেখতে অনেকটা একটা \en{triangle}-এর মতো।}
\sline{আর এই \en{triangle}-টাই হলো —}
\sline{\en{Pascal's Triangle}।}
\clipend
\clip{০৪}{p4\_04}
\sline{কিন্তু এখন প্রশ্ন হলো — এই সংখ্যাগুলো কি শুধু \en{Binomial Theorem} থেকে পাওয়া কিছু \en{coefficient}?}
\sline{নাকি এই \en{triangle}-এর ভেতরেই কোনো নিজস্ব \en{mathematical pattern} লুকিয়ে আছে?}
\sline{চলো, সেটাই দেখি।}
\clipend
\clip{০৫}{p4\_05}
\sline{\en{Triangle}-টার দিকে প্রথমে ভালো করে তাকাও।}
\sline{প্রতিটি \en{row}-এর একদম শুরুতে আছে ওয়ান।}
\sline{আবার একদম শেষেও আছে ওয়ান।}
\sline{যেমন — ওয়ান। ওয়ান, ওয়ান। ওয়ান, টু, ওয়ান। ওয়ান, থ্রি, থ্রি, ওয়ান। ওয়ান, ফোর, সিক্স, ফোর, ওয়ান।}
\sline{এটা খুব সহজেই বোঝা যায়।}
\sline{কারণ প্রতিটি \en{row}-এর দুই প্রান্তের সংখ্যাগুলো হচ্ছে — সি এন জিরো, এবং সি এন এন।}
\sline{আর আমরা জানি — সি এন জিরো সমান ওয়ান, এবং সি এন এন সমান ওয়ান।}
\sline{তাই দুই পাশেই সবসময় ওয়ান থাকবে।}
\clipend
\clip{০৬}{p4\_06}
\sline{ধরো \en{Triangle}-এর এই অংশটা দেখি।}
\sline{এখানে উপরের দুইটা সংখ্যা হলো — ওয়ান, ওয়ান।}
\sline{এদের যোগ করলে পাই — ওয়ান প্লাস ওয়ান সমান টু।}
\sline{আবার একটু নিচে দেখি।}
\sline{এখানে — ওয়ান প্লাস টু সমান থ্রি, এবং টু প্লাস ওয়ান সমান থ্রি।}
\sline{তাই পরের \en{row} হয় — ওয়ান, থ্রি, থ্রি, ওয়ান।}
\sline{আর তার পরেরটাতে — থ্রি প্লাস থ্রি সমান সিক্স।}
\sline{তাই পাই — ওয়ান, ফোর, সিক্স, ফোর, ওয়ান।}
\sline{অর্থাৎ মাঝের প্রতিটি সংখ্যা আসলে তার ঠিক উপরের দুইটি সংখ্যার যোগফল।}
\clipend
\clip{০৭}{p4\_07}
\sline{কিন্তু প্রশ্ন হলো — কেন?}
\sline{এটা কি শুধু একটা সুন্দর \en{pattern}, নাকি এর পেছনে \en{Combination}-এর কোনো কারণ আছে?}
\sline{এখানেই আমাদের আগের \en{Binomial Theorem}-এর \en{concept} আবার কাজে আসে।}
\sline{আমরা জানি, সি এন আর মানে এন টি বস্তুর মধ্যে থেকে আর টি বেছে নেওয়ার সংখ্যা।}
\sline{এখন ধরো এন টি বস্তুর মধ্যে থেকে আর টি বেছে নিচ্ছি।}
\sline{একটা নির্দিষ্ট বস্তু — ধরো এ — নিয়ে চিন্তা করি।}
\sline{আমাদের আর টি বস্তু বেছে নেওয়ার ক্ষেত্রে দুইটি সম্ভাবনা আছে।}
\sline{হয় এ-কে আমরা বেছে নেব। অথবা এ-কে বেছে নেব না।}
\clipend
\clip{০৮}{p4\_08}
\sline{যদি এ-কে বেছে নিই, তাহলে বাকি এন মাইনাস ওয়ান টি বস্তু থেকে আমাদের আরও আর মাইনাস ওয়ান টি নিতে হবে।}
\sline{তাই এই ক্ষেত্রে উপায়ের সংখ্যা — সি, এন মাইনাস ওয়ান, আর মাইনাস ওয়ান।}
\sline{আর যদি এ-কে একদমই না নিই, তাহলে পুরো আর টি বস্তুই বাকি এন মাইনাস ওয়ান টি থেকে নিতে হবে।}
\sline{তাই সেই ক্ষেত্রে উপায়ের সংখ্যা — সি, এন মাইনাস ওয়ান, আর।}
\sline{দুই ক্ষেত্রই একসাথে করলে — সি এন আর সমান, সি এন মাইনাস ওয়ান আর মাইনাস ওয়ান, প্লাস, সি এন মাইনাস ওয়ান আর।}
\clipend
\clip{০৯}{p4\_09}
\sline{এটাই \en{Pascal's Triangle}-এর সবচেয়ে গুরুত্বপূর্ণ \en{pattern}-এর একটি।}
\sline{অর্থাৎ উপরের দুই সংখ্যার যোগফল নিচের সংখ্যায় পরিণত হচ্ছে — বিষয়টা বেশ মজার।}
\sline{কারণ প্রতিবার \en{Combination}-এর ক্যালকুলেশন না করে, সহজেই এই \en{triangle}-এর মাধ্যমে আমরা সহগগুলো বের করতে পারি।}
\sline{এটা সরাসরি \en{Combination}-এর নিয়ম থেকেই এসেছে।}
\clipend
\clip{১০}{p4\_10}
\sline{এবার আরেকটা প্রশ্ন করি।}
\sline{ধরো আমরা একটা সম্পূর্ণ \en{row} নিলাম।}
\sline{যেমন — ওয়ান, থ্রি, থ্রি, ওয়ান।}
\sline{সবগুলো যোগ করলে পাই — ওয়ান প্লাস থ্রি প্লাস থ্রি প্লাস ওয়ান সমান এইট।}
\sline{আরেকটা \en{row} — ওয়ান, ফোর, সিক্স, ফোর, ওয়ান।}
\sline{যোগ করলে — ওয়ান প্লাস ফোর প্লাস সিক্স প্লাস ফোর প্লাস ওয়ান সমান সিক্সটিন।}
\sline{এখানে আবার একটা \en{pattern} দেখা যাচ্ছে।}
\sline{ওয়ান, টু, ফোর, এইট, সিক্সটিন, এভাবে।}
\sline{প্রতিবার আগের \en{row}-এর যোগফলকে টু দিয়ে গুণ করলে পরের \en{row}-এর যোগফল পাওয়া যাচ্ছে।}
\sline{বিষয়টা অনেকটা এরকম।}
\clipend
\clip{১১}{p4\_11}
\sline{আমরা জানি এন-তম \en{row}-এর সংখ্যাগুলো হলো — সি এন জিরো, সি এন ওয়ান, সি এন টু, এভাবে, সি এন এন।}
\sline{তাই পুরো \en{row}-এর যোগফল হলো — সি এন জিরো প্লাস সি এন ওয়ান প্লাস সি এন টু প্লাস, এভাবে, প্লাস সি এন এন।}
\sline{এখন \en{Binomial Theorem}-এর একটি বিশেষ \en{case} চিন্তা করি।}
\sline{যদি এ প্লাস বি হোল টু দ্য পাওয়ার এন-এ, এ সমান ওয়ান এবং বি সমান ওয়ান বসাই, তাহলে বামদিকে পাই — ওয়ান প্লাস ওয়ান হোল টু দ্য পাওয়ার এন, সমান টু টু দ্য পাওয়ার এন।}
\clipend
\clip{১২}{p4\_12}
\sline{আর ডানদিকে \en{Binomial Theorem} অনুযায়ী, এ আর বি-এর মান এক হওয়ায় প্রতিটি \en{term}-এর এ টু দ্য পাওয়ার এন মাইনাস আর, বি টু দ্য পাওয়ার আর অংশটা এক-এ পরিণত হয়।}
\sline{তাই শুধু \en{coefficient}-গুলো পড়ে থাকে — সি এন জিরো প্লাস সি এন ওয়ান প্লাস সি এন টু প্লাস, এভাবে, প্লাস সি এন এন।}
\sline{তাই, সি এন জিরো প্লাস সি এন ওয়ান প্লাস, এভাবে, প্লাস সি এন এন, সমান টু টু দ্য পাওয়ার এন।}
\sline{অর্থাৎ \en{Pascal's Triangle}-এর এন-তম \en{row}-এর সব সংখ্যার যোগফল হলো — টু টু দ্য পাওয়ার এন।}
\clipend
\clip{১৩}{p4\_13}
\sline{এখন \en{visually} দেখি।}
\sline{প্রথম \en{row}: ওয়ান।}
\sline{তারপর — ওয়ান প্লাস ওয়ান সমান টু।}
\sline{তারপর — ওয়ান প্লাস টু প্লাস ওয়ান সমান ফোর।}
\sline{তারপর — ওয়ান প্লাস থ্রি প্লাস থ্রি প্লাস ওয়ান সমান এইট।}
\sline{তারপর — ওয়ান প্লাস ফোর প্লাস সিক্স প্লাস ফোর প্লাস ওয়ান সমান সিক্সটিন।}
\sline{অর্থাৎ — ওয়ান, টু, ফোর, এইট, সিক্সটিন, এভাবে।}
\sline{প্রতিটি ধাপে আগেরটার দ্বিগুণ।}
\sline{আর এটাই হলো — টু টু দ্য পাওয়ার এন।}
\clipend
\clip{১৪}{p4\_14}
\sline{আমরা দেখলাম, \en{Pascal's Triangle}-এর প্রতিটি সংখ্যা আসলে \en{Combination} থেকে এসেছে।}
\sline{দেখলাম, কেন মাঝের সংখ্যা দুটো উপরের সংখ্যার যোগফল হয়।}
\sline{দেখলাম, কেন পুরো একটি \en{row}-এর যোগফল টু টু দ্য পাওয়ার এন হয়।}
\sline{অর্থাৎ শুরুতে যে সংখ্যাগুলোকে আমরা শুধু \en{Binomial Theorem}-এর \en{coefficient} হিসেবে দেখেছিলাম, সেগুলোর নিজেদের মধ্যেও একটা সম্পূর্ণ \en{mathematical structure} রয়েছে।}
\sline{আর এই \en{structure}-টাই হলো — \en{Pascal's Triangle}।}
\clipend
\clip{১৫}{p4\_15}
\sline{তাহলে পুরো বিষয়টাকে একবার শুরু থেকে শেষ পর্যন্ত একসাথে দেখি।}
\sline{প্রথমে আমরা দেখলাম, এন টি বস্তুর মধ্যে থেকে এম টি বেছে নিয়ে দল তৈরি করতে গেলে কেন, এন ইনটু এন মাইনাস ওয়ান ইনটু এন মাইনাস টু, এভাবে, ইনটু এন মাইনাস এম প্লাস ওয়ান, বাই এম ফ্যাক্টোরিয়াল দিয়ে \en{Combination} বের করা যায়।}
\sline{তারপর বুঝলাম, এই \en{Combination}-এর সংখ্যাগুলোই \en{Binomial Theorem}-এর প্রতিটি \en{term}-এর \en{coefficient} তৈরি করে।}
\clipend
\clip{১৬}{p4\_16}
\sline{এরপর এ ওয়ান, এ টু, এ থ্রি এবং বি ওয়ান, বি টু, বি থ্রি-এর মতো আলাদা \en{member} ব্যবহার করে দেখলাম, এ প্লাস বি হোল টু দ্য পাওয়ার এন \en{expand} করার সময় কোন \en{term} কতবার তৈরি হয়।}
\sline{সেখান থেকেই বুঝলাম কেন \en{coefficient} হয় সি এন আর।}
\sline{অর্থাৎ এন টি \en{position}-এর মধ্যে আর টি \en{position}-এ বি নেওয়ার উপায়ের সংখ্যা।}
\sline{তারপর এই \en{coefficient}-গুলোকে আলাদা করে একটার নিচে আরেকটা সাজাতেই আমরা পেলাম — \en{Pascal's Triangle}।}
\clipend
\clip{১৭}{p4\_17}
\sline{এবং এখান থেকে আমরা \en{Pascal's Triangle}-এর দুটি গুরুত্বপূর্ণ \en{pattern} দেখলাম।}
\sline{প্রথমত, উপরের দুই সংখ্যার যোগফল কেন নিচের সংখ্যাটি দেয় — সি এন আর সমান, সি এন মাইনাস ওয়ান আর মাইনাস ওয়ান, প্লাস, সি এন মাইনাস ওয়ান আর।}
\sline{আর দ্বিতীয়ত, একটি সম্পূর্ণ \en{row}-এর সব \en{coefficient} যোগ করলে কেন পাওয়া যায় — টু টু দ্য পাওয়ার এন।}
\clipend
\clip{১৮}{p4\_18}
\sline{মূলত \en{Binomial Theorem} গণিতের খুবই গুরুত্বপূর্ণ একটা অংশ — কারণ এর সাথে \en{calculus}, \en{combination}, \en{statistics} সহ অনেক কিছু সম্পর্কিত।}
\sline{যেমন একটা উদাহরণ হিসেবে বলা যায়, নিউটন যখন তার পাই নির্ণয়ের জন্য যে সূত্রটা তৈরি করেছিলেন, যেটার মাধ্যমে তিনি প্রায় পনেরো ঘর পর্যন্ত পাই-এর মান বের করেছিলেন — সেই \en{infinite series}-টি মূলত এই \en{Binomial Theorem}-এর উপর ভিত্তি করেই তৈরি।}
\clipend
\clip{১৯}{p4\_19}
\sline{যতটুকু পেরেছি চেষ্টা করেছি যে কীভাবে \en{concept}-টা সহজে বোঝানো যায়, অর্থাৎ ভেতরের ধারণা বা \en{logic}-টা কেমন ছিল সেটা তুলে ধরার চেষ্টা করেছি।}
\sline{আজ এ পর্যন্তই। আল্লাহ হাফেজ।}
\clipend
\end{document}
